\documentclass[11pt,a4paper]{article}
\usepackage[margin=25mm]{geometry}
\usepackage{amsmath,amssymb,siunitx,hyperref,listings,microtype}
\hypersetup{colorlinks=true,linkcolor=blue,urlcolor=blue,pdftitle={Report 8 - Grid Model Family for OpenHPL/OpenHPLjl},pdfauthor={Madhusudhan Pandey}}
\lstset{basicstyle=\ttfamily\small,breaklines=true,frame=single}
\title{Report 8 - Grid Model Family for OpenHPL/OpenHPLjl}
\author{Madhusudhan Pandey}
\date{19 September 2026}
\begin{document}
\maketitle
\begin{abstract}
This report introduces the grid model family. The first implementation is an infinite-grid frequency boundary used to close reduced electromechanical models and prepare the package for SMIB studies.
\end{abstract}
\section{Updated Context}
The grid is the downstream electrical boundary:
\[
\boxed{\text{Generator}\rightarrow\text{Grid}}.
\]
For the first SMIB studies, the grid is treated as infinitely strong.

\section{Generalized Concept Model}
An infinite grid fixes electrical frequency while accepting whatever active power is injected by the connected generator.

\section{Other Models Available in the Literature}
A useful hierarchy is
\[
\boxed{\text{infinite bus}\rightarrow\text{frequency-dependent grid}\rightarrow\text{SMIB network}\rightarrow\text{multi-machine grid}}.
\]

\section{Mathematics Involved}
For nominal frequency $f_0$,
\[
\boxed{\omega_g=2\pi f_0}.
\]
At $50~\si{\hertz}$,
\[
\omega_g=100\pi~\si{\radian\per\second}.
\]

\section{OpenHPL-Specific Modeling Interpretation}
OpenHPL contains power-system interfaces intended to couple hydropower plant dynamics to electrical-system models. OpenHPLjl first establishes a reduced active-power/frequency port before introducing voltage-angle network equations.

\section{Implementation in Julia ModelingToolkit / SciML}
\begin{lstlisting}
@connector ElectricalPowerPort begin
    omega(t)
    P(t), [connect = Flow]
end

@component function InfiniteGrid(; name, f0=50.0)
    @named port = ElectricalPowerPort()
    port.omega ~ 2*pi*f0
end
\end{lstlisting}

\section{Model Assembly and Connections}
The ideal generator electrical port connects directly to the infinite grid. The connected generator determines active-power exchange while the grid fixes frequency.

\section{Numerical Experiment / Validation}
For
\[
f_0=50~\si{\hertz},
\]
the expected angular frequency is
\[
\boxed{\omega_g=100\pi\approx314.159~\si{\radian\per\second}}.
\]

\section{Expected Outputs}
Useful outputs are
\[
f(t),\quad\omega(t),\quad P(t).
\]
For the infinite grid, frequency remains constant.

\section{Engineering Interpretation}
The infinite-grid model is a boundary condition, not a dynamic grid. It is appropriate for early SMIB validation before finite grid inertia, load damping, transmission, and inter-area dynamics are introduced.

\section*{Conclusion}
The first electrical boundary is now available. The next milestone is the complete reduced chain
\[
\boxed{\text{Turbine}\rightarrow\text{Shaft}\rightarrow\text{Generator}\rightarrow\text{InfiniteGrid}}
\]
followed by a full hydropower SMIB model.
\end{document}
